% Strongly Solving Okiya: Complete Game-Theoretic Analysis via GPU-Accelerated Brute Force
% arXiv-ready single-file LaTeX paper
\documentclass[11pt,a4paper]{article}

\usepackage[utf8]{inputenc}
\usepackage[T1]{fontenc}
\usepackage{amsmath,amssymb,amsthm}
\usepackage{graphicx}
\usepackage{booktabs}
\usepackage{hyperref}
\usepackage{algorithm}
\usepackage{algpseudocode}
\usepackage{subcaption}
\usepackage[numbers,sort&compress]{natbib}
\usepackage{geometry}
\geometry{margin=1in}

\usepackage{orcidlink}

\hypersetup{
  colorlinks=true,
  linkcolor=blue,
  citecolor=blue,
  urlcolor=blue
}

\newtheorem{theorem}{Theorem}

\title{Strongly Solving Okiya: Complete Game-Theoretic Analysis\\
  and Statistical Win Dynamics via GPU-Accelerated Brute Force}
\author{Aleksei Iurasov\,\orcidlink{0009-0003-5192-8132}}
\date{}

\begin{document}
\maketitle

% ---------------------------------------------------------------------------
% ABSTRACT
% ---------------------------------------------------------------------------
\begin{abstract}
We present the first strong solution of Okiya (Niya), a two-player
combinatorial board game on a $4{\times}4$ grid with attribute-constrained
tile placement and 19 winning patterns.  Using a GPU-accelerated
forward-BFS/backward-minimax solver on an NVIDIA RTX~4090, a single board
configuration is resolved in approximately 20\,s, enumerating
${\sim}9.3{\times}10^{6}$ reachable states---8.7\% of the theoretical
upper bound.  We solve 1{,}000 uniformly random configurations (seed~42,
${\sim}1$\,h total) and find that the first player (P0) wins 67.2\% of
configurations with zero draws, yet the mean number of winning opening moves
is only 1.37 out of~12, revealing a structural but \emph{fragile}
first-mover advantage.  A per-level win-probability gradient $f(L)$
exhibits a pronounced turn-parity sawtooth with shrinking amplitude,
indicating that the game's outcome crystallises progressively rather than at
a single critical level.  Across all configurations, 46.1\% of terminal
states are decided by the stuck rule rather than pattern completion,
highlighting its decisive role in Okiya's dynamics.
\end{abstract}

% ---------------------------------------------------------------------------
% 1. INTRODUCTION
% ---------------------------------------------------------------------------
\section{Introduction}

\subsection{Motivation}

Combinatorial game theory~\cite{berlekamp2018} has produced landmark results
ranging from the exhaustive analysis of Tic-Tac-Toe to the 18-year
computation that solved Checkers~\cite{schaeffer2007}.  The theoretical
foundations---from Zermelo's determinacy theorem for
chess~\cite{zermelo1913} to von Neumann's minimax
theorem~\cite{vonneumann1928}---guarantee that finite two-player
zero-sum games possess optimal strategies, but \emph{finding} those
strategies requires complete enumeration or sophisticated search.
Between the trivially solvable and the computationally intractable lies
a gap: games that are too complex for pencil-and-paper analysis yet
small enough for brute-force computation on modern
hardware~\cite{heule2017}.  Filling this gap is scientifically valuable
because it provides exact benchmarks for heuristic search, validates
game-theoretic predictions, and reveals structural phenomena invisible
to partial analyses.

Okiya---originally published as Niya (2012) and rebranded by Blue Orange
Games as Okiya (2013)~\cite{cathala2013}---is a two-player abstract strategy
game designed by Bruno Cathala.  Played on a $4{\times}4$ grid with
attribute-constrained placement and dual win conditions (pattern completion
and opponent stuck), its state space of ${\sim}9.5{\times}10^{6}$ reachable
positions makes it an ideal candidate for exact solution.  The game is
nontrivial---the relation constraint between successive moves creates a
complex DAG rather than a tree, and the stuck rule introduces a second
strategic axis absent from pure pattern-matching games.

In this work we provide a \emph{strong solution} of individual Okiya
configurations: for each solved configuration, the minimax value is computed
for every reachable state, enabling perfect play from any position.  Because
each of the ${\sim}4.5 \times 10^{9}$ distinct configurations defines a
different game DAG, a strong solution of the full game would require solving
all of them.  We solve 1{,}000 random configurations to characterise the
game's statistical landscape.

\subsection{Related work}

Table~\ref{tab:related} places Okiya in the context of previously solved
combinatorial games.  Retrograde analysis~\cite{thompson1986} and
alpha-beta search~\cite{knuth1975} have been the dominant techniques for
such solutions.  With ${\sim}9.5{\times}10^{6}$ reachable states per
configuration, Okiya sits between Tic-Tac-Toe and Nine Men's Morris in
complexity, occupying a niche where GPU brute force is both sufficient and
informative.

\begin{table}[t]
\centering
\caption{Comparison of solved two-player games (ordered by year solved).
  State-space sizes are approximate and refer to reachable positions.}
\label{tab:related}
\begin{tabular}{@{}llllll@{}}
\toprule
Game & States & Year & Method & Outcome & Ref.\\
\midrule
Tic-Tac-Toe     & $5.5{\times}10^{3}$  & 1952 & Exhaustive     & Draw     & folklore\\
Connect Four     & $4.5{\times}10^{12}$ & 1988 & Alpha-beta     & P1 wins  & \cite{allis1988}\\
Nine Men's Morris& $1.0{\times}10^{10}$ & 1996 & Retrograde     & Draw     & \cite{gasser1996}\\
Checkers         & $5.0{\times}10^{20}$ & 2007 & Retrograde     & Draw     & \cite{schaeffer2007}\\
Pentago          & $3.0{\times}10^{15}$ & 2014 & Alpha-beta+DB  & P1 wins  & \cite{irving2014}\\
\textbf{Okiya}   & $\mathbf{9.5{\times}10^{6}}$\textsuperscript{\textdagger} & \textbf{2026}
  & \textbf{GPU BFS+minimax} & \textbf{P0 67.2\%\textsuperscript{*}} & \textbf{This work}\\
\bottomrule
\multicolumn{6}{@{}l@{}}{\footnotesize\textsuperscript{*}Configuration-dependent; P0 wins 67.2\% of 1{,}000 random configs (95\% CI: [64.3\%, 70.1\%]).}\\
\multicolumn{6}{@{}l@{}}{\footnotesize\textsuperscript{\textdagger}Per configuration; ${\sim}4.5 \times 10^{9}$ distinct configurations exist.}
\end{tabular}
\end{table}

% ---------------------------------------------------------------------------
% 2. GAME RULES AND FORMALIZATION
% ---------------------------------------------------------------------------
\section{Game Rules and Formalization}

\subsection{Board, tiles, and attributes}

Okiya is played on a $4{\times}4$ grid of 16~positions.  Before the game
begins, each position is assigned a unique tile bearing two orthogonal
attributes drawn from:
\begin{itemize}
  \item \textbf{Plant} (4 types): Maple, Cherry, Pine, Iris.
  \item \textbf{Weather} (4 types): Sun, Tanzaku, Bird, Rain.
\end{itemize}
The $4 \times 4 = 16$ attribute combinations yield 16~distinct tile types,
one per position (Figure~\ref{fig:relation}).  A \emph{board configuration}
is a permutation of these 16 tile types on the 16~positions; there are $16!$
possible configurations.

\begin{figure}[t]
  \centering
  \includegraphics[width=0.7\linewidth]{../results/figures/8_relation_graph.png}
  \caption{Tile-relation graph.  Each node is one of the 16~tile types
    (plant${\times}$weather); an edge connects two tiles that share at
    least one attribute.  Every tile has exactly 6~neighbours (degree~6,
    clustering coefficient~0.4), and this structure is invariant across
    all configurations up to relabelling.}
  \label{fig:relation}
\end{figure}

\subsection{Move constraints and win conditions}

Two players, P0 (Red, first mover) and P1 (Blue), alternate turns.  On each
turn a player claims an unclaimed tile subject to the following constraints:
\begin{enumerate}
  \item \textbf{Opening move (turn~0):} P0 must select one of the
    12~\emph{border} positions (all except the four interior squares at
    positions 5, 6, 9, 10 in row-major order).
  \item \textbf{Subsequent moves:} The current player must select an
    unclaimed tile that is \emph{related} to the last-played tile.  Two tiles
    are related if they share at least one attribute (same plant \emph{or}
    same weather).
\end{enumerate}

A player wins by either:
\begin{itemize}
  \item \textbf{Pattern completion:} claiming four tiles that form one of 19
    winning patterns (4~rows, 4~columns, 2~diagonals, 9~overlapping
    $2{\times}2$ blocks); or
  \item \textbf{Stuck rule:} leaving the opponent with no valid move---the
    stuck player loses immediately.
\end{itemize}

\subsection{State encoding}\label{sec:encoding}

Each game state is packed into a single \texttt{uint64} using 37~bits:
\begin{equation}\label{eq:state}
  s = (\mathit{mask}_{P0} \ll 21) \mathbin{|} (\mathit{mask}_{P1} \ll 5) \mathbin{|} \mathit{last\_pos}
\end{equation}
where $\mathit{mask}_{P0}$ and $\mathit{mask}_{P1}$ are 16-bit bitmasks of
each player's claimed positions, and $\mathit{last\_pos} \in \{0,\dots,16\}$
encodes the most recently played position (16~denotes the initial state with
no previous move).  This encoding is injective, sortable (enabling $O(\log
n)$ binary-search lookups), and compact (8~bytes per state).

The \emph{level} of a state is $L = \mathrm{popcount}(\mathit{mask}_{P0}) +
\mathrm{popcount}(\mathit{mask}_{P1})$, the total number of tiles placed.
At level~$L$, P0 has placed $\lceil L/2 \rceil$ tiles and it is P0's turn
if $L$ is even.

\subsection{Configuration space}

The tile arrangement is fixed before play.  There are $16!
\approx 2.09 \times 10^{13}$ possible configurations.  Symmetry reductions
include:
\begin{itemize}
  \item Board symmetry (dihedral group $D_4$, order 8): rotations and
    reflections.
  \item Plant relabelling ($4! = 24$ permutations).
  \item Weather relabelling ($4! = 24$ permutations).
\end{itemize}
Under these reductions, the number of \emph{distinct} configurations is
\[
  \frac{16!}{8 \times 24 \times 24} \approx 4.5 \times 10^{9}.
\]
Each distinct configuration produces a different game DAG with different
state counts, game value, and winning moves.

% ---------------------------------------------------------------------------
% 3. GPU SOLVER ARCHITECTURE
% ---------------------------------------------------------------------------
\section{GPU Solver Architecture}

\subsection{Two-phase overview}

The solver employs a classical \emph{forward enumeration + backward
induction} strategy:
\begin{enumerate}
  \item \textbf{Phase~1 --- Forward BFS:} Starting from the single root
    state (empty board, $L{=}0$), expand all reachable states level by level
    up to $L{=}16$.
  \item \textbf{Phase~2 --- Backward minimax:} Starting from the deepest
    level, propagate game-theoretic values upward to the root.
\end{enumerate}
This produces a \emph{strong solution}: the minimax value and an optimal
strategy are known for every reachable state, not merely the root.

\subsection{Forward BFS}\label{sec:bfs}

Algorithm~\ref{alg:bfs} summarises the forward phase.  At each level, a CUDA
kernel generates all candidate children in parallel (one thread per
parent--position pair, $16n$ threads for $n$~parents), writing valid children
to global memory with a sentinel value (\texttt{0xFFFF\dots F}) for invalid
entries.  The CPU then filters sentinels, deduplicates via
\texttt{np.unique}, detects stuck (no-move) states, and builds a Compressed
Sparse Row (CSR) parent-to-child edge structure for the subsequent minimax
phase.

\begin{algorithm}[t]
\caption{Forward BFS (state enumeration)}\label{alg:bfs}
\begin{algorithmic}[1]
\State $\mathcal{S}_0 \gets \{s_{\text{root}}\}$
\For{$L = 0$ \textbf{to} $15$} \Comment{generates $\mathcal{S}_{L+1}$ up to $\mathcal{S}_{16}$}
  \State \textbf{GPU:} expand non-terminal states in $\mathcal{S}_L$
    \Comment{$16n$ threads}
  \For{each parent $s \in \mathcal{S}_L$, candidate position $p \in \{0,\dots,15\}$}
    \If{$p$ unclaimed \textbf{and} move valid}
      \State $s' \gets \textsc{ApplyMove}(s, p)$
      \If{mover wins} mark $s'$ terminal with value $\pm 1$
      \EndIf
      \State emit $s'$
    \EndIf
  \EndFor
  \State \textbf{CPU:} filter sentinels, deduplicate, detect stuck states
  \State $\mathcal{S}_{L+1} \gets$ unique children
  \State Build CSR edges from $\mathcal{S}_L$ to $\mathcal{S}_{L+1}$
\EndFor
\end{algorithmic}
\end{algorithm}

Move validity at level $L{=}0$ requires a border position ($\mathit{pos}
\mathbin{\&} \text{0xF99F} \neq 0$); at $L{>}0$ it requires the candidate
position to share an attribute with the last-played tile, checked via a
precomputed 16-entry relation-mask table in CUDA constant memory.  Win
detection compares the mover's bitmask against all 19 winning patterns
(Appendix~\ref{app:patterns}).

Deduplication is essential because the game graph is a \emph{DAG}, not a
tree: the same state can be reached by different move orders.  After
\texttt{np.unique}, the sorted state array supports $O(\log n)$ child-index
lookups via \texttt{np.searchsorted}.

\subsection{Backward minimax}\label{sec:minimax}

Algorithm~\ref{alg:minimax} propagates values from terminal states
upward using the classical minimax principle~\cite{knuth1975}.
Values are stored as \texttt{int8}: $+1$ (P0~wins), $-1$ (P1~wins).

\begin{algorithm}[t]
\caption{Backward minimax (value propagation)}\label{alg:minimax}
\begin{algorithmic}[1]
\For{$L = L_{\max} - 1$ \textbf{downto} $0$}
  \For{each non-terminal state $s \in \mathcal{S}_L$} \Comment{GPU, $n$ threads}
    \State $C \gets \text{children of } s$ via CSR
    \If{$L$ even (P0's turn)}
      \State $v(s) \gets \max_{c \in C} v(c)$
        \Comment{exit early on $+1$}
    \Else
      \State $v(s) \gets \min_{c \in C} v(c)$
        \Comment{exit early on $-1$}
    \EndIf
  \EndFor
\EndFor
\end{algorithmic}
\end{algorithm}

Because the value domain is binary ($\{+1, -1\}$, with no draws---see
Section~\ref{sec:nodraw}), the early-exit optimisation is especially
effective: a maximising node can stop on the first $+1$ child and a
minimising node can stop on the first $-1$ child.

\subsection{CUDA implementation details}

Table~\ref{tab:cuda} summarises the GPU resource profile.  The solver is
implemented in PyCUDA~\cite{klockner2012}, following the GPU-based
state-space exploration paradigm of Edelkamp and
Sulewski~\cite{edelkamp2010}.  All layout-dependent data (19~win-pattern
bitmasks, 16~relation masks, 1~border mask) fits in 144~bytes of CUDA
constant memory, residing in L1 cache for the duration of both kernels.
Both kernels use a block size of~256 threads.

\begin{table}[t]
\centering
\caption{GPU memory profile at peak level ($L{=}12$, ${\sim}2.1$\,M parent states).}
\label{tab:cuda}
\begin{tabular}{@{}lrr@{}}
\toprule
Data structure & Per state & Peak total\\
\midrule
States (\texttt{uint64})           & 8\,B   & 17\,MB\\
Values (\texttt{int8})             & 1\,B   & 2\,MB\\
Terminal flags (\texttt{uint8})    & 1\,B   & 2\,MB\\
CSR offsets (\texttt{int32})       & 4\,B   & 9\,MB\\
CSR indices (\texttt{int32})       & var.   & ${\sim}12$\,MB\\
Expansion buffer ($16{\times}8$\,B)& 128\,B & ${\sim}282$\,MB\\
\midrule
Constant memory (win patterns, relation masks, border) & \multicolumn{2}{r}{144\,B}\\
\bottomrule
\end{tabular}
\end{table}

\begin{table}[t]
\centering
\caption{Solver timing breakdown (single configuration, RTX~4090).}
\label{tab:perf}
\begin{tabular}{@{}lr@{}}
\toprule
Phase & Time\\
\midrule
CUDA compilation (cached)     & ${\sim}1$\,s\\
Forward BFS (GPU + CPU dedup) & ${\sim}8$\,s\\
Backward minimax (GPU)        & ${\sim}5$\,s\\
Verification (CPU sampling)   & ${\sim}5$\,s\\
File I/O (compressed npz)     & ${\sim}1$\,s\\
\midrule
\textbf{Total}                & $\mathbf{{\sim}20}$\,\textbf{s}\\
\bottomrule
\end{tabular}
\end{table}

Peak GPU memory is dominated by the expansion buffer at levels 11--12, where
each of ${\sim}2$\,M parent states generates up to 16 candidate children
(${\sim}35$\,M slots, ${\sim}282$\,MB).  The CPU-side bottleneck is
\texttt{np.unique} on arrays of up to ${\sim}35$\,M \texttt{uint64} entries.

Table~\ref{tab:perf} shows the end-to-end timing for a single configuration
on an NVIDIA RTX~4090.  The forward BFS takes approximately 8\,s (GPU
expansion plus CPU deduplication) and backward minimax takes approximately
5\,s.  For the 1{,}000-configuration batch run, CUDA compilation is
amortised and verification is abbreviated, yielding ${\sim}3.5$\,s per
configuration.

\subsection{Verification}

Three independent checks validate the solution:
\begin{enumerate}
  \item \textbf{Structural invariants} (exhaustive): player masks are
    disjoint, popcount matches level, last position belongs to the correct
    player's mask.
  \item \textbf{Minimax consistency} (sampled, up to 50{,}000 states per
    level for interactive runs; exhaustive for the canonical board): each
    non-terminal value equals the max~(P0~turn) or min~(P1~turn) of its
    children's values via CSR lookup.  At ${\sim}9.3$\,M total states,
    exhaustive verification is computationally inexpensive
    (${\sim}2$\,s) and was performed for the canonical configuration.
  \item \textbf{State-count comparison}: per-level counts are printed for
    manual inspection against reference values.
\end{enumerate}

% ---------------------------------------------------------------------------
% 4. CANONICAL BOARD RESULTS
% ---------------------------------------------------------------------------
\section{Canonical Board Results}

We first report detailed results for a single \emph{canonical}
configuration (tile permutation $[14, 4, 11, 9, 7, 12, 10, 3, 8, 1, 5,
0, 6, 2, 13, 15]$ in row-major order, fixed at project inception) before
generalising to 1{,}000 random configurations in
Sections~\ref{sec:statistical}--\ref{sec:dynamics}.

\subsection{Game value}

The canonical board has game value $v(s_{\text{root}}) = +1$: P0 wins with
perfect play.  However, the advantage is extremely fragile: only 1 out of
12 legal opening moves leads to a P0~win.  The remaining 11 opening moves
all result in P1 victories under optimal defence.  At level~1, the fraction
of P0-winning states drops to $f(1) = 1/12 \approx 0.083$, the sharpest
single-move swing in the entire game.

\subsection{State space}

The canonical board contains 9{,}297{,}746 reachable states, approximately
8.7\% of the theoretical upper bound of ${\sim}107$\,M states derived
from combinatorial counting (ignoring the relation constraint).  The
per-level state counts peak at level~12 with 2{,}122{,}834 states (22.8\%
of the total).

Table~\ref{tab:levels} shows the theoretical upper bound and observed
reachable states by level.  The relation constraint prunes approximately
91\% of theoretically possible states, with pruning most severe at early
levels and relaxing toward the endgame.

\begin{table}[t]
\centering
\caption{Theoretical upper bound vs.\ reachable states by level (canonical board).}
\label{tab:levels}
\begin{tabular}{@{}rrrrr@{}}
\toprule
Level & Upper bound & Reachable & \% & Mover\\
\midrule
0  & 1           & 1           & 100.0 & P0\\
1  & 16          & 12          &  75.0 & P1\\
2  & 480         & 72          &  15.0 & P0\\
3  & 5{,}040     & 360         &   7.1 & P1\\
4  & 43{,}680    & 1{,}524     &   3.5 & P0\\
5  & 218{,}400   & 6{,}866     &   3.1 & P1\\
6  & 960{,}960   & 24{,}198    &   2.5 & P0\\
7  & 2{,}802{,}800 & 90{,}332  &   3.2 & P1\\
8  & 7{,}207{,}200 & 244{,}503 &   3.4 & P0\\
9  & 12{,}972{,}960 & 671{,}170 &  5.2 & P1\\
10 & 20{,}180{,}160 & 1{,}227{,}894 & 6.1 & P0\\
11 & 22{,}198{,}176 & 2{,}101{,}611 & 9.5 & P1\\
12 & 20{,}180{,}160 & 2{,}122{,}834 & 10.5 & P0\\
13 & 12{,}492{,}480 & 1{,}772{,}951 & 14.2 & P1\\
14 & 5{,}765{,}760  & 779{,}608    & 13.5 & P0\\
15 & 1{,}544{,}400  & 229{,}834    & 14.9 & P1\\
16 & 205{,}920      & 23{,}976     & 11.6 & P0\\
\midrule
\textbf{Total} & \textbf{106{,}778{,}593} & \textbf{9{,}297{,}746} & \textbf{8.7} &\\
\bottomrule
\end{tabular}
\end{table}

\subsection{Terminal state classification}

Of the 9{,}297{,}746 reachable states, 1{,}938{,}231 (20.8\%) are terminal.
These divide into two categories:
\begin{itemize}
  \item \textbf{Pattern wins:} 1{,}083{,}641 states (55.9\% of terminals)
    where a player completed one of the 19 winning patterns.
  \item \textbf{Stuck wins:} 854{,}590 states (44.1\% of terminals) where
    the active player has no legal move and therefore loses.
\end{itemize}
The stuck rule is nearly as impactful as pattern completion in determining
game outcomes.  The earliest terminal states appear at level~7 (P0's fourth
tile, the minimum for a four-in-a-pattern win); stuck states first appear at
level~8.  The mean game length is 12.66 moves, with the mode at level~13
(30.0\% of terminal states).

\subsection{Win rate under random play}

Beyond minimax values, we compute the probability of a P0~win under
\emph{uniform random play} (each player selects uniformly among valid moves)
via bottom-up aggregation over the full DAG.  From the root, P0's win
probability under random play is \textbf{51.9\%}---a modest advantage
compared to the decisive $+1$ under optimal play.  This divergence
illustrates that many of P1's minimax-losing positions require precise
defensive play to exploit, making P0's advantage practically less dominant
than the game value suggests.

% ---------------------------------------------------------------------------
% 5. STATISTICAL ANALYSIS (1000 CONFIGS)
% ---------------------------------------------------------------------------
\section{Statistical Analysis: 1{,}000 Configurations}\label{sec:statistical}

\subsection{Methodology}

We sample 1{,}000 board configurations by permuting the 16 tile types using
a seeded pseudo-random generator (seed~42, NumPy default RNG).  These
configurations are sampled from the raw $16!$ permutation space without
symmetry deduplication; given the ${\sim}4.5 \times 10^{9}$ equivalence
classes, the probability of any two samples being symmetry-equivalent is
negligible ($< 10^{-4}$ by the birthday bound).  Each
configuration is solved end-to-end (BFS + minimax) with per-level metrics
extracted in-memory.  The total computation takes approximately 1~hour on a
single RTX~4090 (${\sim}3.5$\,s per configuration), with CUDA compilation
amortised across runs.

\subsection{Outcome distribution}

Figure~\ref{fig:outcomes} shows the outcome distribution.  Of 1{,}000
configurations:
\begin{itemize}
  \item \textbf{672 (67.2\%, 95\% Wald CI: [64.3\%, 70.1\%])} are won by P0
    (first player).
  \item \textbf{328 (32.8\%)} are won by P1 (second player).
  \item \textbf{0 (0.0\%)} end in draws.
\end{itemize}
The absence of draws across ${\sim}9.48 \times 10^{9}$ aggregate reachable
states is consistent with the structural draw impossibility proven in
Theorem~1 (Section~\ref{sec:nodraw}).

\begin{figure}[t]
  \centering
  \includegraphics[width=0.8\linewidth]{../results/figures/1_win_distribution.png}
  \caption{Outcome distribution across 1{,}000 random configurations.
    P0 wins 672 configurations (67.2\%), P1 wins 328 (32.8\%), with zero
    draws.}
  \label{fig:outcomes}
\end{figure}

\subsection{State-space variation}

The total number of reachable states is remarkably stable across
configurations: mean 9{,}481{,}800, standard deviation 204{,}000 (2.2\%
coefficient of variation), range 8{,}731{,}257--10{,}120{,}199.  This
near-invariance stems from the uniform relation graph: every tile is related
to exactly 6 others (degree~6.0, clustering coefficient~0.4, identical
across all configurations up to relabelling).  The combinatorial structure of
the game graph is thus dominated by the attribute-relation topology rather
than the specific tile placement.

\subsection{Winning openings}\label{sec:openings}

Figure~\ref{fig:openings} shows the distribution of winning opening moves
(out of 12~border positions) for P0.  The mean is 1.37 (95\% CI:
[1.26, 1.48]) with a median of~1 and a mode of~0.
\begin{itemize}
  \item 328 configurations (32.8\%) have \emph{zero} winning openings (P1
    wins).
  \item 292 (29.2\%) have exactly 1 winning opening.
  \item 196 (19.6\%) have 2 winning openings.
  \item 184 (18.4\%) have 3--8 winning openings.
\end{itemize}
Even when P0 wins, the advantage is narrow: nearly half of P0-winning
configurations (292 of 672, 43.5\%) offer only a single correct first move.

\begin{figure}[t]
  \centering
  \includegraphics[width=0.8\linewidth]{../results/figures/5_winning_openings.png}
  \caption{Histogram of the number of winning opening moves for P0 across
    1{,}000 configurations.  The mode is~0 (P1-winning configurations);
    among P0 wins, a single winning opening is most common.}
  \label{fig:openings}
\end{figure}

% ---------------------------------------------------------------------------
% 6. WIN DYNAMICS
% ---------------------------------------------------------------------------
\section{Win Dynamics}\label{sec:dynamics}

\subsection{Win-probability gradient}

For each configuration, define the \emph{P0-win fraction} at level~$L$ as
\begin{equation}\label{eq:fL}
  f(L) = \frac{|\{s \in \mathcal{S}_L : v(s) = +1\}|}{|\mathcal{S}_L|},
\end{equation}
the fraction of reachable states at level~$L$ whose minimax value is a P0
win.  Since Okiya is zero-sum with no draws, the P1-win fraction is
$g(L) = 1 - f(L)$.

Figure~\ref{fig:gradient} plots the mean $f(L)$ across 1{,}000
configurations with $\pm 1$ standard deviation bands.  The curve does not
rise monotonically; instead it exhibits a pronounced \emph{sawtooth}
superimposed on a gradually converging trend.  The confidence bands are tight
(std $< 0.05$ for $L \geq 6$), confirming that the sawtooth is a universal
property of Okiya, not an artefact of specific configurations.

\begin{figure}[t]
  \centering
  \includegraphics[width=0.85\linewidth]{../results/figures/2_win_probability_gradient.png}
  \caption{P0-win fraction $f(L)$ across game levels, averaged over 1{,}000
    configurations ($\pm 1$ std shaded).  The sawtooth pattern reflects
    turn-parity coupling: $f(L)$ rises on even levels (P0's turn) and falls
    on odd levels (P1's turn).}
  \label{fig:gradient}
\end{figure}

\subsection{Turn-parity sawtooth}

The sawtooth arises from the coupling of level and turn.  At even levels, P0
moves and can steer toward favourable positions ($f$ increases); at odd
levels, P1 moves and pushes back ($f$ decreases).  Figure~\ref{fig:sawtooth}
isolates this effect by plotting the first difference $\Delta f(L) = f(L) -
f(L{-}1)$ as a bar chart with error bars across configurations.

The sawtooth amplitude is largest at the opening ($|\Delta f(1)| \approx
0.56$, reflecting the drop from $f(0) = 0.672$---which equals the
aggregate P0-win rate, since level~0 contains exactly one state per
configuration---to $f(1) = 0.11$) and
decreases toward mid-game ($|\Delta f| \approx 0.2$--$0.3$ at levels
8--12), before a final spike at levels 15--16 driven by the stuck rule
endgame.  The shrinking amplitude indicates that individual moves carry less
leverage as the game progresses---the outcome is increasingly determined by
the accumulated position rather than any single move.

\begin{figure}[t]
  \centering
  \includegraphics[width=0.85\linewidth]{../results/figures/3_sawtooth.png}
  \caption{First difference $\Delta f(L) = f(L) - f(L{-}1)$ of the P0-win
    fraction.  Positive bars (even levels) indicate P0's turn improving its
    position; negative bars (odd levels) indicate P1's countermoves.  Error
    bars show $\pm 1$ std across 1{,}000 configs.}
  \label{fig:sawtooth}
\end{figure}

\subsection{Critical move ratio}

For each non-terminal state, we check whether any child has a different
minimax value than the parent.  The \emph{critical move ratio} (CMR) at
level~$L$ is the fraction of non-terminal states at that level where at
least one such value-changing child exists---i.e., where the current
player's choice is decisive.

On the canonical board, $\mathrm{CMR}(0) = 0.92$: eleven of twelve opening
moves change the game's outcome.  By mid-game ($L \approx 6$--$10$),
$\mathrm{CMR} \approx 0.3$, and by the endgame ($L \geq 14$), $\mathrm{CMR}
\to 0$---the outcome is already determined regardless of remaining moves.
Figure~\ref{fig:cmr} shows this monotonic decline across all 1{,}000
configurations.

\begin{figure}[t]
  \centering
  \includegraphics[width=0.85\linewidth]{../results/figures/6_critical_move_ratio.png}
  \caption{Critical move ratio (CMR) by level, averaged over 1{,}000
    configurations ($\pm 1$ std shaded).  The ratio declines from
    ${\sim}0.9$ at the opening to near zero by level~14, quantifying
    how early the game's outcome crystallises.}
  \label{fig:cmr}
\end{figure}

\subsection{Stuck-state fraction}

The fraction of terminal states decided by the stuck rule (rather than
pattern completion) grows monotonically from 0\% at level~7 to 76\% at
level~16 across the 1{,}000-configuration aggregate.  Early terminations
($L \leq 9$) are dominated by pattern wins, while late terminations ($L \geq
14$) are overwhelmingly stuck outcomes.  Overall, 46.1\% of all terminal states (across
${\sim}1.93 \times 10^{9}$ aggregate terminals) are stuck wins.  The
per-configuration mean stuck fraction is 46.1\% with a 95\% CI of
[45.3\%, 46.9\%] over the 1{,}000 configurations, indicating low
cross-configuration variance.  The stuck rule is thus a co-equal
determinant of Okiya's outcome.
Figure~\ref{fig:stuck} illustrates this level-by-level
transition from pattern-dominated to stuck-dominated terminations.

\begin{figure}[t]
  \centering
  \includegraphics[width=0.85\linewidth]{../results/figures/7_stuck_fraction.png}
  \caption{Fraction of terminal states decided by the stuck rule
    (vs.\ pattern completion) by level, aggregated over 1{,}000
    configurations.  Stuck outcomes dominate from level~14 onward,
    reaching 76\% at level~16.}
  \label{fig:stuck}
\end{figure}

% ---------------------------------------------------------------------------
% 7. HEURISTIC DISTILLATION
% ---------------------------------------------------------------------------
\section{Heuristic Distillation}\label{sec:heuristic}

A strong solution is a lookup table with millions of entries---useful for
theoretical analysis but impractical as a cognitive strategy.  A natural
follow-up question is: how much of the minimax solution can be
\emph{distilled} into concise, human-playable rules?  We address this by
extracting interpretable heuristics directly from the labeled game DAG of
the canonical configuration and evaluating them against the oracle.

\subsection{Feature engineering and mutual information}

We define 11~features per (state, candidate move) pair, spanning four
categories: \emph{offensive} (pattern completion, threat count),
\emph{defensive} (opponent pattern progress, block urgency),
\emph{constraint} (opponent mobility after the move, stuck threat), and
\emph{positional} (pattern participation count, game level).

To identify which features matter at each game phase, we compute the
mutual information $\mathrm{MI}(\text{feature}, \text{is\_optimal})$
stratified by level on a sample of ${\sim}23{,}000$ informative states
(those with both optimal and suboptimal moves, yielding ${\sim}80{,}000$
labeled move pairs).  Table~\ref{tab:mi} summarises the results.

\begin{table}[t]
\centering
\caption{Dominant features by game phase (mutual information with
  optimal-move label).  The strategic axis shifts from positional play
  (early) through pattern building (mid) to constraint play (late).}
\label{tab:mi}
\begin{tabular}{@{}lll@{}}
\toprule
Level range & Top features (MI) & Interpretation\\
\midrule
$L{=}0$    & None (MI${=}0$ for all) & Opening is a memorisation problem\\
$L{=}1$--3 & Pattern participation (0.25) & Positional value dominates\\
$L{=}4$--6 & Own best pattern (0.07) & Pattern building\\
$L{=}7$--9 & Opponent mobility (0.08) & Constraint play emerges\\
$L{=}10$--12 & Own best pattern (0.18), & Three-way race:\\
             & opp.\ mobility (0.17), win-now (0.16) & win, constrain, or lose\\
$L{=}13$--14 & Win immediately (0.69) & Near-deterministic endgame\\
\bottomrule
\end{tabular}
\end{table}

The phase transition from positional features to constraint features
independently confirms the structural observation from
Section~\ref{sec:dynamics}: Okiya's strategy operates on two axes
(pattern completion and stuck threat), with the balance shifting from
the former to the latter as the game progresses.

\subsection{Heuristic hierarchy and move accuracy}

We evaluate a family of heuristics ranging from trivial to learned,
measuring \emph{move accuracy}: the fraction of informative states where
the heuristic selects a minimax-optimal move.
Table~\ref{tab:heuristics} presents the results.

\begin{table}[t]
\centering
\caption{Heuristic hierarchy: move accuracy (agreement with minimax
  oracle) and win rate against a uniform-random opponent (500~games,
  averaged over both roles).  Canonical configuration.}
\label{tab:heuristics}
\begin{tabular}{@{}llrr@{}}
\toprule
ID & Heuristic & Accuracy & Win rate\\
\midrule
H0 & Random & 0.41 & 50.0\%\\
H1 & Win-or-block (2~rules) & 0.52 & 65.6\%\\
H2 & + Minimise opp.\ mobility (3~rules) & 0.62 & 74.1\%\\
H3 & + Maximise own patterns (4~rules) & 0.65 & 82.6\%\\
H4 & Decision tree, depth~4 (${\sim}5$ splits) & 0.74 & 83.7\%\\
H5 & Decision tree, depth~8 (${\sim}30$ splits) & 0.78 & 84.5\%\\
\midrule
Oracle & Minimax lookup & 1.00 & 99.8\%\textsuperscript{$\ddagger$}\\
\bottomrule
\multicolumn{4}{@{}l@{}}{\footnotesize\textsuperscript{$\ddagger$}The
  Oracle loses when playing as the losing side (P1 on a P0-winning board);
  even perfect play cannot overcome a lost position.}
\end{tabular}
\end{table}

A four-rule heuristic (\emph{win if you can; block if you must; build
on high-value positions; constrain the opponent}) closes roughly
one-third of the gap between random play and perfection.  A shallow
decision tree (depth~4) closes 56\%\ of the gap, discovering
level-dependent thresholds that hand-crafted rules miss.

The depth-4 tree yields five human-readable rules:
\begin{enumerate}
  \item If you can complete a win pattern, play it.
  \item If the opponent will be stuck (0~valid moves), play it.
  \item If the opponent will have only 1~valid move, play it.
  \item At high-value positions (${\geq}6$ pattern participations):
    play if the opponent has no imminent 3-in-a-pattern threat;
    otherwise block.
  \item Otherwise, avoid the move.
\end{enumerate}

\noindent In natural language: \emph{Win if you can.\ \ Trap if you
can't win.\ \ Build on valuable positions, but watch for threats.}

\subsection{The opening bottleneck}\label{sec:opening-bottleneck}

A striking result is that \emph{every} heuristic achieves 0\% accuracy
at $L{=}0$.  On the canonical board, all 12~border positions produce
identical feature vectors: each participates in exactly 4~win patterns,
each has 6~related tiles, and no blocking or stuck threats exist on an
empty board.  The winning opening depends on the \emph{global} structure
of the tile-relation graph---which positions create advantageous
constraint cascades 5--10~moves later---not on any locally computable
property.

This confirms quantitatively that the ``fragile first-mover advantage''
identified in Sections~\ref{sec:openings} and~\ref{sec:dynamics} is
concentrated in board-specific global structure rather than locally
computable features.  A practical Okiya strategy would therefore combine
\emph{opening preparation} (memorise the winning first move for the
given board) with \emph{heuristic play} from $L{=}1$ onward---paralleling
chess, where opening theory is memorised and heuristic evaluation governs
the middle game.

\subsection{Limitations of the heuristic analysis}

The heuristic evaluation is conducted on the canonical board
configuration only.  While the near-invariant state-space size across
1{,}000 configurations (2.2\% CV, Section~\ref{sec:statistical})
suggests that the MI profile and accuracy numbers should generalise,
this remains to be verified empirically.  Multi-configuration heuristic
evaluation is left for future work.

% ---------------------------------------------------------------------------
% 8. DISCUSSION
% ---------------------------------------------------------------------------
\section{Discussion}

\subsection{Classification}

Following the taxonomy of van den Herik et al.~\cite{vandenherik2002}, our
result constitutes a \emph{strong solution} of Okiya: the minimax value is
computed for every reachable state of each solved configuration.  In
complexity, Okiya sits between Tic-Tac-Toe ($5.5 \times 10^{3}$ states,
trivially solved) and Nine Men's Morris ($10^{10}$ states, requiring
retrograde analysis~\cite{gasser1996}).  The GPU brute-force approach is both
sufficient and informative at this scale.

\subsection{Draw impossibility}\label{sec:nodraw}

We prove that draws are impossible in Okiya.

\begin{theorem}
Every reachable Okiya state has minimax value $v(s) \in \{+1, -1\}$.
No draw ($v(s) = 0$) can occur under any configuration.
\end{theorem}

\begin{proof}
We proceed by backward induction on the game DAG.

\emph{Base case (terminal states).}  A state $s$ is terminal if either
(a)~the current player's move completes one of the 19~winning patterns,
in which case $v(s) = +1$ or $-1$ depending on the winner; or (b)~the
current player has no legal move (stuck), in which case that player
loses, giving $v(s) \in \{+1, -1\}$.  There is no third terminal condition: at level~16, either the
16th tile itself completes a winning pattern (a pattern win at the
final move) or no unclaimed tiles remain and the current player is
trivially stuck.  In both cases the outcome is decisive.
Thus every terminal state has value in $\{+1, -1\}$.

\emph{Inductive step (non-terminal states).}  Let $s$ be a
non-terminal state at level~$L$.  By minimax, $v(s) = \max_{c \in C}
v(c)$ if it is P0's turn, or $v(s) = \min_{c \in C} v(c)$ if it is
P1's turn, where $C$ is the set of children of~$s$.  By the inductive
hypothesis, $v(c) \in \{+1, -1\}$ for all $c \in C$.  Since
$\max\{+1, -1\} = +1$, $\min\{+1, -1\} = -1$, $\max\{+1\} = +1$,
$\min\{-1\} = -1$, etc., the result $v(s) \in \{+1, -1\}$ follows in
all cases.
\end{proof}

This structural property is confirmed empirically: across 1{,}000
configurations and ${\sim}9.48 \times 10^{9}$ aggregate reachable states,
zero draws were observed.

\subsection{First-mover advantage: structural but fragile}

P0 wins 67.2\% of random configurations---a significant first-mover
advantage.  Yet the advantage is \emph{fragile}: the mean number of winning
opening moves is only 1.37 out of~12, and 43.5\% of P0-winning
configurations offer only a single correct opening.  Under uniform random
play (51.9\% P0 win rate on the canonical board), the advantage nearly
vanishes, indicating that it requires precise play to realise.

The fragility is captured by the P0-win fraction at level~1: $f(1) = 0.11$
averaged across configurations, meaning that after P0's first move, only
11\% of resulting positions are P0-winning.  The ``advantage'' is better
described as \emph{access to a narrow winning path} rather than a broad
strategic superiority.

\subsection{Generalizability}

The two-phase solver (forward BFS + backward minimax on a DAG) is applicable
to any game satisfying:
\begin{enumerate}
  \item \textbf{Bounded state space:} all reachable states fit in GPU memory
    (${\sim}10^{7}$--$10^{8}$ states with current hardware).
  \item \textbf{DAG structure:} the game graph is acyclic, enabling
    level-by-level processing.
  \item \textbf{Sortable encoding:} states can be packed into fixed-width
    integers for efficient deduplication and lookup.
\end{enumerate}
Candidate games include Quarto variants (attribute-based placement on a grid),
modified placement games with relation constraints, and $N{\times}N$
generalisations of Okiya.

\subsection{Strategic compressibility}

The heuristic analysis in Section~\ref{sec:heuristic} suggests a
preliminary observation about game complexity.  Define the
\emph{heuristic ceiling} $\kappa(k)$ as the maximum move accuracy
achievable by a depth-$k$ decision tree trained on the complete
solution: $\kappa(k) = \max_{T \in \mathcal{T}_k} \mathrm{Acc}(T)$,
where $\mathcal{T}_k$ is the set of all depth-$k$ trees over the
feature space.  For Okiya (canonical configuration), $\kappa(4) = 0.74$
and $\kappa(8) = 0.78$, with an estimated asymptote of
${\sim}0.85$--$0.90$.  The ${\sim}15$--20\% residual gap represents
decisions that require deep lookahead or board-specific global knowledge
and resist compression into local features.  Two games with similar
state-space sizes might exhibit very different heuristic ceilings,
indicating different \emph{strategic} complexities.  We note this as a
potentially useful complement to state-space size, though validation
across multiple games and configurations is needed before it can be
considered a robust complexity measure.

\subsection{Limitations}

This work has several limitations.  First, only 1{,}000 of
${\sim}4.5 \times 10^{9}$ distinct configurations are solved; while the
low variance in state-space size and the confidence intervals on the
P0-win rate suggest the sample is representative, configuration-level
features that predict whether a board is P0- or P1-winning remain
unexplored.  Second, the heuristic analysis
(Section~\ref{sec:heuristic}) is evaluated on the canonical
configuration only; multi-configuration validation is needed.  Third, the
solver is benchmarked on a single GPU (NVIDIA RTX~4090, 24\,GB VRAM);
performance on smaller GPUs is not characterised, though peak memory
usage (${\sim}282$\,MB) suggests broad compatibility.  Fourth, the
1{,}000-configuration batch uses sampled (not exhaustive) minimax
verification; exhaustive verification was performed only for the
canonical configuration.  Finally, we do not compare against
alternative solving strategies (e.g., depth-first alpha-beta), which
may be faster for determining the root value alone without computing
all state values.

% ---------------------------------------------------------------------------
% 8. CONCLUSION
% ---------------------------------------------------------------------------
\section{Conclusion and Open Questions}\label{sec:open}

\paragraph{Contributions.}
This work makes four contributions:
\begin{enumerate}
  \item \textbf{Strong solution of individual Okiya configurations:} the
    first complete game-theoretic analysis, computing minimax values for
    all ${\sim}9.3 \times 10^{6}$ reachable states per configuration in
    ${\sim}20$\,s on an RTX~4090.
  \item \textbf{Statistical landscape:} a 1{,}000-configuration study
    showing P0 wins 67.2\% with zero draws, a fragile advantage (mean 1.37
    winning openings), and near-invariant state-space size
    (std/mean~$= 2.2$\%).
  \item \textbf{Win-dynamics analysis:} identification of the turn-parity
    sawtooth in $f(L)$, the stuck-rule dominance at late levels (76\% at
    $L{=}16$), and the shrinking critical move ratio quantifying strategic
    crystallisation.
  \item \textbf{Heuristic distillation:} extraction of interpretable
    playing rules from the complete solution, showing that a 4-rule
    heuristic achieves 65\% move accuracy and a depth-4 decision tree
    reaches 74\% (depth-8: 78\%), while the opening ($L{=}0$) resists
    all feature-based
    approximation---confirming the fragile first-mover advantage is
    rooted in global board structure.
\end{enumerate}

\paragraph{Open questions.}
Several directions remain:
\begin{itemize}
  \item \textbf{GNN outcome prediction:} can a graph neural network on the
    tile-relation graph predict the game value without solving?
  \item \textbf{Full enumeration:} solving all ${\sim}4.5 \times 10^{9}$
    distinct configurations ($4.5 \times 10^{9} \times 3.5\,\text{s}
    \approx 500$~GPU-years on a single RTX~4090; tractable with
    large-scale parallelism).
  \item \textbf{Human play study:} the heuristic hierarchy
    (Section~\ref{sec:heuristic}) brackets human-playable performance
    between 41\% (random) and 78\% (learned tree) move accuracy.
    Empirical studies with human participants would locate actual play
    on this spectrum and test whether the depth-4 rules improve
    human performance.
  \item \textbf{Multi-configuration heuristic validation:} do the MI
    profiles and move-accuracy results from the canonical board
    generalise across all 1{,}000 solved configurations?
  \item \textbf{$N{\times}N$ complexity:} extending Okiya's mechanics to
    larger grids---what are the complexity thresholds for tractability?
\end{itemize}

% ---------------------------------------------------------------------------
% APPENDICES
% ---------------------------------------------------------------------------
\appendix

\section{Win Pattern Bitmasks}\label{app:patterns}

The 19~winning patterns are encoded as 16-bit bitmasks over the $4{\times}4$
board (row-major order, position $p$ corresponds to bit $2^{p}$).  A player
wins if $(\mathit{mask} \mathbin{\&} \mathit{pattern}) = \mathit{pattern}$
for any pattern.

\begin{table}[ht]
\centering
\caption{Win pattern bitmasks.}
\begin{tabular}{@{}llr@{}}
\toprule
Category & Hex & Decimal\\
\midrule
Row 0       & \texttt{0xF000} & 61440\\
Row 1       & \texttt{0x0F00} & 3840\\
Row 2       & \texttt{0x00F0} & 240\\
Row 3       & \texttt{0x000F} & 15\\
\midrule
Column 0    & \texttt{0x8888} & 34952\\
Column 1    & \texttt{0x4444} & 17476\\
Column 2    & \texttt{0x2222} & 8738\\
Column 3    & \texttt{0x1111} & 4369\\
\midrule
Diagonal $\backslash$ & \texttt{0x8421} & 33825\\
Diagonal $/$ & \texttt{0x1248} & 4680\\
\midrule
Block (0,0) & \texttt{0xCC00} & 52224\\
Block (0,1) & \texttt{0x6600} & 26112\\
Block (0,2) & \texttt{0x3300} & 13056\\
Block (1,0) & \texttt{0x0CC0} & 3264\\
Block (1,1) & \texttt{0x0660} & 1632\\
Block (1,2) & \texttt{0x0330} & 816\\
Block (2,0) & \texttt{0x00CC} & 204\\
Block (2,1) & \texttt{0x0066} & 102\\
Block (2,2) & \texttt{0x0033} & 51\\
\bottomrule
\end{tabular}
\end{table}

\section{Notation}\label{app:notation}

\begin{table}[ht]
\centering
\caption{Notation reference.}
\begin{tabular}{@{}ll@{}}
\toprule
Symbol & Meaning\\
\midrule
P0, P1 & Player 0 (Red, first mover), Player 1 (Blue)\\
$L$ & Level: total number of tiles placed\\
$\mathit{mask}_{P0}$, $\mathit{mask}_{P1}$ & 16-bit bitmasks of claimed positions\\
$\mathit{last\_pos}$ & Position of most recently placed tile ($0$--$15$, or $16{=}$none)\\
$v(s)$ & Minimax value of state $s$: $+1$ or $-1$\\
$f(L)$ & P0-win fraction at level $L$ (Eq.~\ref{eq:fL})\\
$\Delta f(L)$ & First difference $f(L) - f(L{-}1)$\\
CMR & Critical move ratio\\
CSR & Compressed Sparse Row (graph storage format)\\
DAG & Directed acyclic graph\\
BFS & Breadth-first search\\
\bottomrule
\end{tabular}
\end{table}

% ---------------------------------------------------------------------------
% CODE AVAILABILITY
% ---------------------------------------------------------------------------
\section*{Code Availability}

Source code, solver, batch-analysis scripts, and all data used in this paper
are available at
\url{https://github.com/format37/okiya}.

% ---------------------------------------------------------------------------
% BIBLIOGRAPHY
% ---------------------------------------------------------------------------
\begin{thebibliography}{99}

\bibitem{zermelo1913}
E. Zermelo,
``\"Uber eine Anwendung der Mengenlehre auf die Theorie des Schachspiels,''
in \emph{Proc.\ Fifth Int.\ Congress of Mathematicians}, vol.~2,
pp.~501--504, 1913.

\bibitem{vonneumann1928}
J. von Neumann,
``Zur Theorie der Gesellschaftsspiele,''
\emph{Mathematische Annalen}, vol.~100, no.~1, pp.~295--320, 1928.

\bibitem{knuth1975}
D.~E. Knuth and R.~W. Moore,
``An analysis of alpha-beta pruning,''
\emph{Artificial Intelligence}, vol.~6, no.~4, pp.~293--326, 1975.

\bibitem{thompson1986}
K. Thompson,
``Retrograde analysis of certain endgames,''
\emph{ICCA Journal}, vol.~9, no.~3, pp.~131--139, 1986.

\bibitem{allis1988}
V. Allis,
``A knowledge-based approach of Connect-Four,''
Master's thesis, Vrije Universiteit Amsterdam, 1988.

\bibitem{gasser1996}
R. Gasser,
``Solving Nine Men's Morris,''
\emph{Computational Intelligence}, vol.~12, no.~1, pp.~24--41, 1996.

\bibitem{vandenherik2002}
H.~J. van den Herik, J.~W.~H.~M. Uiterwijk, and J. van Rijswijck,
``Games solved: Now and in the future,''
\emph{Artificial Intelligence}, vol.~134, no.~1--2, pp.~277--311, 2002.

\bibitem{schaeffer2007}
J. Schaeffer, N. Burch, Y. Bj\"ornsson, A. Kishimoto, M. M\"uller,
R. Lake, P. Lu, and S. Sutphen,
``Checkers is solved,''
\emph{Science}, vol.~317, no.~5844, pp.~1518--1522, 2007.

\bibitem{edelkamp2010}
S. Edelkamp and D. Sulewski,
``Efficient explicit-state model checking on general purpose graphics
processors,''
in \emph{Proc.\ SPIN}, LNCS 6349, pp.~106--123, 2010.

\bibitem{klockner2012}
A. Kl\"ockner, N. Pinto, Y. Lee, B. Catanzaro, P. Ivanov, and
A. Fasih,
``PyCUDA and PyOpenCL: A scripting-based approach to GPU run-time code
generation,''
\emph{Parallel Computing}, vol.~38, no.~3, pp.~157--174, 2012.

\bibitem{cathala2013}
B. Cathala,
\emph{Okiya / Niya},
Blue Orange Games, 2013.
BoardGameGeek ID: 125311.

\bibitem{irving2014}
G. Irving, J. Donkers, and J. Uiterwijk,
``Solving Pentago,''
Technical report, 2014.

\bibitem{berlekamp2018}
E.~R. Berlekamp, J.~H. Conway, and R.~K. Guy,
\emph{Winning Ways for Your Mathematical Plays}, 2nd ed.
A~K Peters/CRC Press, 2018.

\bibitem{heule2017}
M.~J.~H. Heule and O. Kullmann,
``The science of brute force,''
\emph{Communications of the ACM}, vol.~60, no.~8, pp.~70--79, 2017.

\end{thebibliography}

\end{document}
